\chapter{Implementación y resultados}
\label{ch:implementacion}

Este capítulo detalla el proceso de construcción del simulador: las decisiones técnicas concretas tomadas en cada componente, los problemas que surgieron durante el desarrollo y la solución adoptada en cada caso, y los resultados del entrenamiento y evaluación de los modelos de elección modal.

El sistema implementado está disponible como código abierto en \url{https://github.com/ivanuclm/movilidad-urbana}, y permite su reproducción, extensión y reutilización en futuros trabajos de investigación o divulgación sobre movilidad urbana. El procedimiento de despliegue para reproducir el entorno local se detalla en el Anexo~\ref{anx:despliegue}.


% -------------------------------------------------------
\section{Infraestructura de enrutado viario: OSRM}
\label{sec:impl_osrm}
% -------------------------------------------------------

\subsection{Origen y características del extracto OSM}
\label{subsec:impl_osrm_origen}

Open Source Routing Machine (OSRM) trabaja sobre un extracto de OpenStreetMap (OSM), un fichero binario con extensión \texttt{.osm.pbf} que contiene la geometría de calles, carreteras y caminos de un territorio junto con sus atributos de velocidad, sentido y accesibilidad por modo de transporte. Geofabrik \cite{GeofabrikDownloads} distribuye estos extractos por jerarquía geográfica, desde continentes hasta comunidades autónomas.

Se utilizó el extracto de Castilla-La Mancha al ser la región en la que se enmarca el proyecto, con Toledo como caso de estudio, tal como se detalla en la Sección~\ref{sec:impl_otp}. OpenTripPlanner reutiliza el mismo extracto para la red viaria peatonal, de modo que un único fichero sirve a dos servicios del sistema. La Tabla~\ref{tab:osm_extractos} resume los extractos de OSM evaluados durante el desarrollo.

\begin{table}[H]
  \caption{Comparativa de extractos OSM disponibles en Geofabrik.}
  \label{tab:osm_extractos}
  \centering
  \begin{tabular}{lrl}
    \toprule
    \textbf{Extracto} & \textbf{Tamaño (.osm.pbf)} & \textbf{Uso en el proyecto} \\
    \midrule
    Castilla-La Mancha  & \textasciitilde{}98 MB     & Producción (OSRM y OTP) \\
    % Comunidad Valenciana & \textasciitilde{}131 MB     & Pruebas iniciales con feed EMT Valencia \\
    Comunidad de Madrid & \textasciitilde{}79 MB     & Pruebas con feed EMT Madrid \\
    España              & \textasciitilde{}1{,}3 GB  & Evaluado, descartado por tamaño \\
    Europa              & \textasciitilde{}32{,}0 GB & Descartado por tamaño \\
    \bottomrule
  \end{tabular}
\end{table}

\subsection{Evolución del despliegue}
\label{subsec:impl_osrm_evolucion}

El despliegue de OSRM pasó por tres iteraciones antes de alcanzar la configuración definitiva.

\textbf{Fase 1: servidor de demostración oficial.} La primera implementación utilizó el servidor de demostración público del proyecto OSRM (\texttt{router.project-osrm.org}) \cite{OSRMDocs}. Este servidor ofrece acceso gratuito a la API de enrutado sin necesidad de infraestructura local y está pensado para pruebas rápidas de integración. El problema es que solo procesa el perfil de conducción independientemente del indicado en la URL \cite{OSRMIssue4034}; cambiar el perfil en la petición no tiene efecto, de modo que la comparación multiperfil que el simulador requiere resultaba inviable.

\textbf{Fase 2: instancias públicas FOSSGIS.} Las instancias públicas mantenidas por FOSSGIS \cite{FOSSGISRouting} sí devuelven rutas diferenciadas por perfil mediante endpoints independientes (\texttt{/routed-car}, \texttt{/routed-bike}, \texttt{/routed-foot}), lo que permitió confirmar que el esquema de tres perfiles independientes funcionaba y ajustar la integración del backend. El inconveniente es que aplican limitaciones de uso (\textit{rate limiting}), presentan latencia variable y no permiten fijar la versión de los datos, comprometiendo así la reproducibilidad.

\textbf{Fase 3: instancia local con Docker.} La solución definitiva fue desplegar OSRM localmente mediante la imagen oficial \texttt{osrm/osrm-backend} \cite{OSRMGitHub}, con un contenedor independiente por modo de transporte (conducción, bicicleta y a pie). OSRM solo admite un perfil por proceso: cada instancia se preprocesa con un fichero de reglas distinto y sirve únicamente ese modo, de ahí la necesidad de tres contenedores independientes. El transporte público en autobús es responsabilidad de OpenTripPlanner. Esta arquitectura elimina la dependencia de servicios externos en tiempo de ejecución, asegura reproducibilidad y suprime los límites de cuota. El preprocesado de los grafos, que originalmente se ejecutaba a mano, se automatizó como un paso de inicialización de la orquestación Docker Compose (Sección~\ref{sec:arch_infra}), de modo que el despliegue completo no requiere ningún comando manual adicional.


\subsection{Pipeline de preprocesado por perfil}
\label{subsec:impl_osrm_pipeline}

Para que OSRM pueda calcular rutas, el extracto \texttt{.osm.pbf} debe transformarse en una estructura de datos interna optimizada para ejecutar consultas y obtener el camino mínimo entre dos puntos del mapa. Este preprocesado se ejecuta una sola vez por perfil y solo hay que repetirlo si se actualiza el extracto de OSM.

Cada perfil define las reglas de movilidad del modo correspondiente, esto es, qué tipos de vía son accesibles, a qué velocidad y con qué restricciones de giro. Los perfiles se escriben en Lua, el lenguaje de script utilizado por OSRM para describir las reglas de cada modo de transporte. Los tres ficheros de perfil estándar (\texttt{car.lua}, \texttt{bicycle.lua} y \texttt{foot.lua}) están incluidos en la imagen oficial de Docker \texttt{osrm/osrm-backend} en \texttt{/opt/}, por lo que no es necesario descargarlos ni configurarlos externamente. Todas las etapas del preprocesado se ejecutan con esa misma imagen, asegurando la reproducibilidad en cualquier entorno de ejecución.

\begin{enumerate}
  \item \textbf{Extracción} (\texttt{osrm-extract}): lee el \texttt{.osm.pbf} y construye el grafo viario interno aplicando las reglas del perfil Lua indicado: qué vías son transitables, a qué velocidad y con qué restricciones. El resultado es el fichero \texttt{.osrm} junto con varios ficheros auxiliares de índice.
  \item \textbf{Particionado} (\texttt{osrm-partition}): divide el grafo en celdas jerárquicas para el algoritmo Multi-Level Dijkstra (MLD) \cite{Geisberger2008CH}. Esta partición es la que permite calcular rutas de larga distancia en milisegundos, al limitar la búsqueda a fronteras entre celdas en lugar de explorar el grafo completo.
  \item \textbf{Personalización} (\texttt{osrm-customize}): asigna los pesos finales de cada arista sobre la estructura de celdas anterior, produciendo el grafo listo para servir peticiones de ruta.
\end{enumerate}

Las tres etapas generan los ficheros \texttt{clm.osrm} y sus auxiliares, almacenados en directorios separados por perfil (\texttt{osrm-clm/car/}, \texttt{osrm-clm/bike/}, \texttt{osrm-clm/foot/}). Por su tamaño, estos artefactos se versionan mediante \textit{Git LFS} (véase la Sección~\ref{sec:arch_infra}), de modo que un clon del repositorio dispone de los grafos sin necesidad de regenerarlos.

El preprocesado completo de los tres perfiles sobre el extracto de Castilla-La Mancha tarda \todo{Revisar tiempos} aproximadamente 97 segundos en la máquina de desarrollo. Los artefactos generados ocupan en total unos 2,1~GB: el perfil de conducción produce 291~MB porque solo incorpora la red viaria accesible al tráfico rodado, mientras que los perfiles de bicicleta y peatonal alcanzan los 911~MB cada uno al incluir además senderos, caminos rurales y zonas de acceso restringido al tráfico.

\begin{figure}[htbp]
  \centering
  % \missingfigure{Diagrama del pipeline de preprocesado OSRM: \texttt{.osm.pbf} $\to$ \texttt{osrm-extract -p \{perfil\}.lua} $\to$ \texttt{osrm-partition} $\to$ \texttt{osrm-customize} $\to$ \texttt{osrm-routed -{}-algorithm mld}. Indicar que el proceso se repite de forma independiente para los tres perfiles (car, bike, foot), generando tres directorios de artefactos separados.}
  \includegraphics[width=1.0\textwidth]{figs/Pipeline_Preprocesado_OSRM.pdf}

  \caption{Pipeline de preprocesado OSRM por perfil de transporte.}
  \label{fig:osrm_pipeline}
\end{figure}

\subsection{Despliegue y verificación}
\label{subsec:impl_osrm_despliegue}

Cada uno de los tres contenedores OSRM ejecuta el servidor de enrutado sobre su grafo preprocesado con el algoritmo MLD (Figura~\ref{fig:osrm_pipeline}):
\begin{lstlisting}[language=bash]
osrm-routed --algorithm mld /data/clm.osrm
\end{lstlisting}

Los tres contenedores comparten imagen y comando de arranque. Cada uno monta el directorio de artefactos de su perfil (\texttt{osrm-clm/car}, \texttt{osrm-clm/bike}, \texttt{osrm-clm/foot}). Dentro de la red interna de Docker Compose, el proceso \texttt{osrm-routed} escucha siempre en el puerto 5000 en todos los contenedores, y el backend los referencia por nombre de servicio (\texttt{osrm-car}, \texttt{osrm-bike}, \texttt{osrm-foot}) en ese mismo puerto. Para verificar las instancias directamente desde el terminal del anfitrión, cada contenedor expone un puerto distinto: 5001, 5002 y 5003, respectivamente para los perfiles de conducción, bicicleta y a pie. El puerto 5000 queda libre en el anfitrión porque en macOS lo tiene reservado por defecto el receptor AirPlay; esta publicación de puertos se usa solo para depuración, ya que el backend se comunica con OSRM por la red interna.

Una respuesta válida devuelve \texttt{"code": "Ok"} junto con la distancia en metros, la duración en segundos y la geometría en formato de polilínea codificada~\cite{GooglePolyline}: una cadena de texto que comprime la secuencia de coordenadas de la ruta mediante diferencias relativas y codificación Base64, que el backend decodifica antes de enviar las coordenadas al frontend. Si el grafo está incompleto o el par origen-destino cae parcialmente fuera del extracto de Castilla-La Mancha, OSRM devuelve \texttt{"code": "NoRoute"} o una ruta incompleta, detectable al comparar visualmente los tres trazados en la interfaz.

\begin{code}[H]
\begin{lstlisting}[language=bash]
# osrm-car  (conduccion, localhost:5001)
curl "http://localhost:5001/route/v1/driving/-4.0356,39.8750;-4.0023,39.8602?overview=full"
# osrm-bike (bicicleta,  localhost:5002)
curl "http://localhost:5002/route/v1/cycling/-4.0356,39.8750;-4.0023,39.8602?overview=full"
# osrm-foot (peatonal,   localhost:5003)
curl "http://localhost:5003/route/v1/foot/-4.0356,39.8750;-4.0023,39.8602?overview=full"
\end{lstlisting}
\caption{Verificación de las tres instancias OSRM desde el terminal del anfitrión. El parámetro \texttt{overview=full} fuerza la devolución de la geometría completa de la ruta; sin él, OSRM devuelve una versión simplificada con menos puntos.}
\label{cod:osrm_verificacion}
\end{code}

La Figura~\ref{fig:osrm_rutas} muestra las tres rutas calculadas simultáneamente para el mismo par origen-destino: en azul, la ruta de conducción sigue las calles principales; en verde, la ruta de bicicleta prioriza carriles bici y calles tranquilas; en gris con trazado discontinuo, la ruta peatonal accede a zonas restringidas al tráfico rodado. Las diferencias entre los tres trazados reflejan las distintas restricciones de cada perfil sobre la red viaria de Toledo.

\begin{figure}[H]
  \centering
  \includegraphics[width=1.0\textwidth]{figs/RutasSoloOSRM.jpeg}
  \caption{Rutas OSRM para los tres perfiles de transporte sobre Toledo.}
  \label{fig:osrm_rutas}
\end{figure}

% -------------------------------------------------------
\section{Planificación multimodal: OpenTripPlanner}
\label{sec:impl_otp}
% -------------------------------------------------------

\subsection{Fuente de datos GTFS y proceso de selección}
\label{subsec:impl_otp_gtfs}

OpenTripPlanner requiere dos fuentes de datos: un extracto viario OSM para modelar los tramos de acceso a pie a las paradas, y un feed GTFS con el servicio programado del operador de transporte público. La fuente utilizada para los feeds fue el Punto de Acceso Nacional de Transporte Multimodal (NAP), gestionado por el Ministerio de Transportes y Movilidad Sostenible, que centraliza feeds GTFS de operadores de toda España \cite{NAPHome}. El extracto OSM de Castilla-La Mancha ya disponible para OSRM (Sección~\ref{subsec:impl_osrm_origen}) se reutilizó directamente, sin necesidad de preparar ningún fichero adicional.

La disponibilidad de feeds de transporte urbano en Castilla-La Mancha resultó muy limitada. En Albacete solo figuraban servicios interurbanos como Alsa, sin datos de red urbana en el NAP. Guadalajara contaba con feeds del Consorcio de Transportes de Madrid, pero sin red municipal propia. Cuenca tampoco disponía de feed urbano. La única opción urbana era Ciudad Real, cuyo feed incluía también Seseña, con paradas repartidas entre dos localidades separadas por más de 170~km; además, en ese momento no incluía geometría de rutas, requisito imprescindible para la representación cartográfica.

Ante la ausencia de un feed válido en Castilla-La Mancha, se evaluaron feeds de otras ciudades. La EMT de Madrid \cite{NAPEMTMadrid}, con 4.914 paradas, 236 rutas y 87.137 viajes programados, sirvió de banco de pruebas para validar la integración de OTP con datos reales y desarrollar los endpoints de la capa estática del backend. El volumen del feed reveló un problema de paginación: la API exige un límite explícito y devolvía solo las primeras 500 paradas por defecto, dejando la red visualmente incompleta. El límite se ajustó a 5.000, suficiente para Madrid y los feeds posteriores.

El feed GTFS elegido finalmente fue el de los autobuses urbanos de Toledo \cite{NAPToledoUrbano}. Con 268 paradas, 56 rutas y 87.751 viajes que cubren aproximadamente tres meses de servicio programado, tiene un tamaño ajustado al alcance del simulador, abarca el casco urbano, el polígono industrial y las urbanizaciones periféricas, e incluye geometría de rutas para su representación cartográfica. Como el feed no incluye tarifas, el coste del billete se introduce como parámetro en el panel de la interfaz, junto al resto de variables del perfil de viajero, para su uso en la inferencia modal.

\subsection{Periodo de validez del feed y selección de fecha}
\label{subsec:impl_otp_fecha}

Los feeds GTFS del NAP tienen un periodo de validez limitado; el del operador de Toledo (\texttt{GTFS\_Urbano\_Toledo\_2026.zip}) cubre del 22 de febrero al 22 de mayo de 2026. Si la fecha de consulta cae fuera de ese rango, OTP no devuelve itinerarios con tramos de transporte público y el simulador trata la situación como si no hubiera servicio disponible, devolviendo únicamente un trayecto a pie equivalente al que proporciona OSRM con el perfil peatonal. Las particularidades del feed de Toledo relativas a la estructura de sus líneas se describen en la Sección~\ref{sec:impl_gtfs}.

El mismo feed alimenta los dos subsistemas que lo consumen: el grafo de OTP (Sección~\ref{subsec:impl_otp_grafo}) y la capa estática de horarios (Sección~\ref{sec:impl_gtfs}), de modo que ambos comparten exactamente la misma ventana de validez. La fecha y la hora de consulta se exponen mediante un selector global en la interfaz, con un valor por defecto dentro del rango del feed (\texttt{2026-05-21}, \texttt{12:00}) y acotado al intervalo válido. Este selector es transversal a toda la simulación: gobierna a la vez los itinerarios multimodales, los horarios de la red y el día y la hora del perfil de inferencia, evitando configuraciones inconsistentes entre las distintas vistas. La ubicación y el comportamiento de este selector en la interfaz se describen en la Sección~\ref{subsec:impl_frontend_mapa}. El backend recibe la fecha y la hora en la petición; cuando no se indican, aplica el valor por defecto.

Se fijó la hora por defecto a las 12:00 por situarse en la franja de mayor frecuencia de servicio y evitar el horario nocturno, donde la oferta es reducida o inexistente.

\subsection{Construcción del grafo multimodal}
\label{subsec:impl_otp_grafo}

El grafo se construye con un único comando Docker ejecutado sobre el directorio \texttt{otp-toledo/}, que debe contener el extracto OSM (\texttt{clm.osm.pbf}) y el feed GTFS (\texttt{GTFS\_Urbano\_Toledo\_2026.zip}):

\begin{lstlisting}[language=bash]
docker run --rm -v "otp-toledo:/var/opentripplanner" opentripplanner/opentripplanner:2.5.0 --build --save
\end{lstlisting}

OTP lee ambos ficheros y serializa el resultado en \texttt{graph.obj}. La red viaria OSM modela los tramos de acceso a pie a las paradas, y el feed GTFS aporta el servicio programado con horarios, paradas y geometría de las líneas. El proceso tarda unos minutos \todo{tiempos} y solo es necesario repetirlo cuando se actualiza el feed; el extracto OSM puede reutilizarse entre versiones. El comando anterior corresponde a la generación manual del grafo.

El \texttt{graph.obj} resultante (unos 120\,MB) se versiona mediante \textit{Git LFS} (véase la Sección~\ref{sec:arch_infra}), de modo que el sistema arranca sin reconstrucción en un clon limpio. Como mecanismo de reserva, la orquestación incluye un servicio de inicialización idempotente (\texttt{otp-build}) que ejecuta \texttt{-{}-build -{}-save} únicamente si el grafo no está presente, de forma similar a la compilación de los grafos OSRM. Una vez disponible, OTP se lanza en modo de servicio (\texttt{-{}-load -{}-serve}) y queda accesible en el puerto 8080, exponiendo la API REST de planificación de rutas \cite{OTPDocs}.



\subsection{Integración de itinerarios en el backend}
\label{subsec:impl_otp_integracion}

El router \texttt{/api/otp} del backend consulta OTP a través del endpoint REST de planificación \texttt{/otp/routers/default/plan}, que es el punto de acceso estándar de OTP~2.x para solicitar itinerarios entre dos coordenadas. OTP devuelve hasta cinco itinerarios alternativos. El criterio de selección automática es el siguiente: se elige el primero que incluya al menos un tramo de transporte público (modo distinto de \texttt{WALK}); si ninguno cumple esta condición, se devuelve el de menor duración total. El frontend permite paginar entre las alternativas mediante los botones ``Anterior'' y ``Siguiente'', que transmiten el índice del itinerario deseado al backend, de modo que el vector de características para la inferencia se calcula siempre sobre el mismo itinerario que el usuario visualiza.

Para cada tramo del itinerario seleccionado se devuelve: modo de transporte, distancia, duración, geometría decodificada desde el formato de polilínea codificada, parada de inicio y fin, nombre de línea, agencia y horarios de salida y llegada en formato \texttt{HH:MM}. Para los tramos de transporte público se añade a la consulta el parámetro \texttt{showIntermediateStops}, que incorpora a la respuesta de OTP las paradas intermedias de cada tramo; con ellas el backend compone la secuencia ordenada completa de paradas del tramo, desde el embarque hasta el desembarque, anotando para cada una su nombre, coordenadas y hora de paso. Esta secuencia alimenta el diagrama de paradas del itinerario descrito en la Sección~\ref{subsec:impl_frontend_rutas}. La geometría de cada tramo se combina para componer la traza completa del itinerario, que el frontend renderiza diferenciando los segmentos por modo de transporte.

OTP devuelve los instantes de paso en tiempo universal (epoch UTC) acompañados del desfase horario de la agencia (\texttt{agencyTimeZoneOffset}). El backend suma ese desfase antes de formatear las horas, de modo que los itinerarios se muestran en hora local de Toledo (Europe/Madrid), con el horario de verano resuelto automáticamente según la fecha consultada.

Cuando OTP devuelve un itinerario compuesto exclusivamente por un tramo a pie, las características de transporte público del vector de entrada tomarían valores próximos a cero, que en el dataset caracterizan precisamente un transporte público rápido y directo. La solución adoptada, descrita en la Sección~\ref{subsec:impl_backend_penalizacion}, consiste en sobrescribir esas características con valores de penalización extremos que llevan al modelo a descartar el transporte público en ausencia de servicio real.

La Figura~\ref{fig:otp_itinerario} en la Sección~\ref{subsec:impl_frontend_rutas} ilustra un itinerario multimodal típico en la interfaz del simulador.

% -------------------------------------------------------
\section{Integración del feed GTFS}
\label{sec:impl_gtfs}
% -------------------------------------------------------

El formato GTFS (\textit{General Transit Feed Specification}) \cite{GTFSOverview} estructura la información de transporte público como una colección de ficheros de texto plano con extensión \texttt{.txt} y formato de valores separados por comas, empaquetados en un único archivo ZIP. En el arranque del servicio, el backend descomprime y carga en memoria con \texttt{pandas} los seis ficheros necesarios: \texttt{stops.txt} (paradas con coordenadas y nombre), \texttt{routes.txt} (líneas de la red), \texttt{trips.txt} (servicios programados por línea), \texttt{stop\_times.txt} (horarios de paso por cada parada), \texttt{shapes.txt} (geometría de los trazados) y \texttt{calendar\_dates.txt} (calendario de excepciones de servicio) \cite{GTFSReference}. Los \textit{DataFrames} resultantes permanecen en memoria durante toda la vida del proceso, evitando releer el feed en cada petición y eliminando la necesidad de un sistema gestor de base de datos para información que no cambia en tiempo de ejecución.

Responder a preguntas como <<¿qué líneas pasan por esta parada?>> o <<¿cuál es el horario de una ruta para una fecha concreta?>> requiere cruzar varios de estos ficheros, tal como ilustra la Figura~\ref{fig:gtfs_joins}. Para obtener las líneas asociadas a una parada se encadenan tres cruces: \texttt{stops.txt} $\to$ \texttt{stop\_times.txt} $\to$ \texttt{trips.txt} $\to$ \texttt{routes.txt}, unidos por \texttt{stop\_id}, \texttt{trip\_id} y \texttt{route\_id}, respectivamente. Para los horarios por fecha se añade el filtro de \texttt{calendar\_dates.txt}, que indica los días en que cada servicio está activo.

\begin{figure}[htbp]
  \centering
  \missingfigure{Diagrama de los seis ficheros GTFS como cajas (\texttt{stops}, \texttt{stop\_times}, \texttt{trips}, \texttt{routes}, \texttt{shapes}, \texttt{calendar\_dates}) con flechas que muestran los cruces por clave: \texttt{stops}\,--\,\texttt{stop\_times} (stop\_id), \texttt{stop\_times}\,--\,\texttt{trips} (trip\_id), \texttt{trips}\,--\,\texttt{routes} (route\_id), \texttt{trips}\,--\,\texttt{shapes} (shape\_id) y \texttt{trips}\,--\,\texttt{calendar\_dates} (service\_id). Resaltar la cadena stops$\to$stop\_times$\to$trips$\to$routes que resuelve ``líneas por parada''.}
  \caption{Cruces entre los ficheros del feed GTFS para resolver las consultas de la capa estática.}
  \label{fig:gtfs_joins}
\end{figure}

El feed de Toledo presenta una particularidad en la estructura de sus líneas. La empresa operadora, identificada en los feeds del NAP como UNAUTO S.L. (Grupo Ruiz), registra cada sentido de recorrido como un \texttt{route\_id} independiente con el mismo \texttt{route\_short\_name}: el feed contiene 56 entradas en \texttt{routes.txt} que corresponden a solo 25 líneas distintas. Por ejemplo, la línea L5 tiene asignados los identificadores \texttt{50011} y \texttt{50012}, uno por cada sentido de circulación. El backend agrupa los \texttt{route\_id} con el mismo \texttt{route\_short\_name} en un diccionario de hermanos (\texttt{siblingRouteIds}), de modo que el catálogo de líneas muestra una sola entrada por línea. Al seleccionar una parada concreta desde el mapa, que conoce el \texttt{route\_id} exacto del servicio que pasa por ella, se activa directamente el sentido correcto sin lógica adicional. Los endpoints que expone el router GTFS a partir de esta estructura se describen en la Sección~\ref{subsec:impl_backend_gtfs}.

% -------------------------------------------------------
\section{Implementación del backend FastAPI}
\label{sec:impl_backend}
% -------------------------------------------------------

\subsection{Estructura modular}
\label{subsec:impl_backend_estructura}

El backend está implementado con FastAPI sobre Python y organizado en módulos por responsabilidad. Un directorio \texttt{api/} contiene los cuatro routers (\texttt{routes\_osrm.py}, \texttt{routes\_otp.py}, \texttt{routes\_gtfs.py}, \texttt{routes\_lpmc.py}); un directorio \texttt{services/} contiene la lógica de negocio desacoplada de la capa HTTP (\texttt{osrm\_client.py}, \texttt{lpmc\_inference.py}); y el módulo raíz \texttt{main.py} monta los routers y habilita CORS (\textit{Cross-Origin Resource Sharing}), el mecanismo que autoriza al navegador a que la interfaz web, servida desde un puerto distinto, consuma la API del backend. Los endpoints expuestos por cada router se recogen en la Tabla~\ref{tab:endpoints_api} del capítulo~\ref{ch:arquitectura}.

La separación en módulos permite desarrollar, probar y documentar cada router por separado. Además, FastAPI genera de forma automática documentación interactiva siguiendo el estándar OpenAPI en la ruta \texttt{/docs}. Durante el desarrollo, esta documentación permitió lanzar peticiones a cada endpoint y comprobar sus respuestas directamente desde el navegador, sin escribir código adicional de prueba.

\begin{figure}[htbp]
  \centering
  \missingfigure{Captura de la documentación OpenAPI automática (\texttt{/docs}) del backend FastAPI, mostrando los cuatro grupos de endpoints (/api/osrm, /api/otp, /api/gtfs, /api/lpmc) con sus rutas, métodos HTTP y esquemas de petición y respuesta.}
  \caption{Documentación OpenAPI generada automáticamente por FastAPI.}
  \label{fig:fastapi_docs}
\end{figure}

\subsection{Router OSRM: consultas multiperfil}
\label{subsec:impl_backend_osrm}

El router \texttt{/api/osrm} expone el endpoint \texttt{POST /api/osrm/routes}, que acepta un par origen-destino y la lista de perfiles de transporte solicitados. Las consultas a los tres contenedores OSRM se lanzan en paralelo mediante la primitiva \texttt{asyncio.gather}, de modo que el tiempo total de respuesta equivale al de la consulta más lenta y no a la suma de las tres.

Para cada perfil, OSRM devuelve la distancia en metros, la duración en segundos y la geometría en formato de polilínea codificada. El backend decodifica esa geometría en una lista de coordenadas \texttt{[lat, lon]} para que pueda enviarse al frontend y representarse con Leaflet. Por último, las duraciones se convierten de segundos a horas antes de construir el vector de características, ya que el dataset LPMC almacena todas las duraciones en esa unidad.

\subsection{Router OTP: itinerarios multimodales}
\label{subsec:impl_backend_otp}

El router \texttt{/api/otp} expone el endpoint \texttt{POST /api/otp/routes}, que consulta el planificador de OTP, gestiona la selección de itinerario (automática o según el itinerario que el usuario tenga seleccionado en el frontend) y transforma la respuesta de OTP a un formato propio y homogéneo de tramos, más sencillo de consumir por el frontend. Para cada tramo del itinerario seleccionado se devuelven el modo de transporte, la distancia, la duración, la geometría decodificada desde la polilínea codificada, las paradas de inicio y fin, el nombre de la línea, la agencia y los horarios de salida y llegada. Las geometrías de los tramos se combinan en una sola traza del itinerario, que el frontend vuelve a separar para dibujar cada tramo según su modo de transporte.

\subsection{Router GTFS: capa estática}
\label{subsec:impl_backend_gtfs}

El router \texttt{/api/gtfs} expone la información estática de la red de autobús directamente a partir de los ficheros del feed GTFS cargados en memoria (Sección~\ref{sec:impl_gtfs}), sin depender de OpenTripPlanner. Esto permite consultar y visualizar la red aunque el planificador no esté disponible. Expone cuatro endpoints:

\begin{itemize}
  \item \textbf{Listado de paradas} (\texttt{GET /api/gtfs/stops}): devuelve todas las paradas del feed con identificador, nombre, coordenadas y líneas asociadas. Acepta un parámetro \texttt{limit} (5.000 por defecto, suficiente para el feed de Toledo con 268 paradas) y un filtro de caja delimitadora opcional para restringir el resultado al área visible en el mapa.

  \item \textbf{Listado de líneas} (\texttt{GET /api/gtfs/routes}): catálogo completo de líneas con identificador, nombre corto, nombre largo y atributos de color cuando están disponibles en el feed.

  \item \textbf{Detalle de línea} (\texttt{GET /api/gtfs/routes/\{route\_id\}}): dado el identificador de una línea, devuelve la secuencia ordenada de paradas y la geometría del trazado asociado para su representación cartográfica.

  \item \textbf{Horarios por fecha} (\texttt{GET /api/gtfs/routes/\{route\_id\}/schedule}): dado un identificador de ruta y una fecha mediante el parámetro \texttt{?date=YYYY-MM-DD}, devuelve los servicios programados agrupados por sentido, con primera y última salida y listado completo de horarios. Al igual que con OTP (Sección~\ref{subsec:impl_otp_fecha}), la interfaz usa una fecha dentro del rango de validez del feed; cualquier fecha posterior devuelve cero servicios.
\end{itemize}

La lógica de agrupación de sentidos descrita en la Sección~\ref{sec:impl_gtfs} (los 56 \texttt{route\_id} del feed de Toledo agrupados en 25 líneas por \texttt{route\_short\_name}) se aplica tanto en este router como en el frontend, de modo que el comportamiento es coherente en todo el sistema.

\subsection{Penalización de itinerarios sin transporte público real}
\label{subsec:impl_backend_penalizacion}

Cuando OTP devuelve un itinerario compuesto exclusivamente por un tramo a pie, sin ningún tramo de transporte público, las características de transporte público del vector de entrada (tiempo de acceso a la parada, tiempo en el vehículo, tiempo de espera y de transbordo, y número de transbordos) tomarían valores próximos a cero. El problema es que esos mismos valores bajos caracterizan en el dataset a un transporte público rápido y directo, por lo que el modelo podría asignar una probabilidad elevada al transporte público en un trayecto donde en realidad no hay servicio. El caso se detectó durante las pruebas de integración del módulo de inferencia: en trayectos cortos en los que la opción más rápida era ir a pie, OTP devolvía ese único tramo como mejor itinerario y la inferencia resultaba engañosa.

Un problema detectado reside en que los modelos de aprendizaje automático que serán entrenados en la Sección~\ref{sec:impl_lpmc} realizan la inferencia con un número fijo y predefinido de modos de transporte. De esta forma, si el viaje no admite transporte público, no es posible eliminar este modo para la inferencia. Para solucionar este problema, el backend comprueba si el itinerario contiene algún tramo de transporte público. Si no es así, sobrescribe las seis características de transporte público (\texttt{dur\_pt\_access}, \texttt{dur\_pt\_bus}, \texttt{dur\_pt\_rail}, \texttt{dur\_pt\_int\_waiting}, \texttt{dur\_pt\_int\_walking} y \texttt{pt\_n\_interchanges}) con un valor situado en el extremo superior de su distribución. En concreto, a cada característica se le asigna el valor que, tras el escalado estándar (\texttt{StandardScaler}) del modelo, queda cinco desviaciones típicas por encima de la media de entrenamiento: es decir, $\mu + 5\sigma$, tomando la media $\mu$ y la desviación típica $\sigma$ que el escalador aprendió de esa característica. De este modo el modelo percibe un transporte público con tiempos y transbordos extremos, prácticamente inviable, y descarta esa opción sin necesidad de modificar su arquitectura ni el número de variables de entrada.

De forma análoga, en trayectos muy cortos (menos de 500~m de distancia en línea recta, calculada con la fórmula de Haversine) el backend añade un recargo fijo de 10 minutos al tiempo de conducción del coche, que representa el aparcamiento y el acceso al vehículo. Sin él, el modelo tendería a elegir el coche siempre que su tiempo de conducción fuese menor, ignorando que para distancias tan cortas el sobrecoste de aparcar hace más razonable ir a pie. Ambas correcciones son ejemplos de conocimiento experto inyectado en las entradas del modelo; el análisis de cómo responde cada modelo a la penalización y la justificación de este enfoque se abordan en la sección~\ref{subsec:impl_lpmc_entrenamiento}.

\subsection{Router LPMC: inferencia modal y descubrimiento de modelos}
\label{subsec:impl_backend_lpmc}

El router \texttt{/api/lpmc} centraliza la lógica de inferencia de elección modal y el acceso a los modelos entrenados. A diferencia de los routers de enrutado, que delegan el cálculo en servicios externos (OSRM y OTP), este router ejecuta la predicción de forma local sobre los artefactos serializados en \texttt{lpmc/models/}. Expone cuatro endpoints:

\begin{itemize}
  \item \textbf{Predicción} (\texttt{POST /api/lpmc/predict}). Recibe el par origen-destino, el itinerario de OTP que el usuario tiene seleccionado, las variables sociodemográficas del viajero y, de forma opcional, el nombre del modelo a usar (\texttt{model\_variant}). Cuando no se indica el nombre del modelo, se recurre a la variable de entorno \texttt{LPMC\_MODEL\_VARIANT}. La respuesta contiene la probabilidad de cada uno de los cuatro modos de transporte (a pie, bicicleta, transporte público y coche), el modo más probable para ese viajero junto con su probabilidad, y el vector de características de ruta con el que se ha inferido. Incluye además algunos metadatos útiles para depuración, entre ellos el indicador \texttt{pt\_available}, que vale verdadero cuando el itinerario contiene al menos un tramo de transporte público y falso cuando no hay servicio real (caso descrito en la Sección~\ref{subsec:impl_backend_penalizacion}).

  \item \textbf{Comparación de modelos} (\texttt{POST /api/lpmc/compare}). Ejecuta la predicción con los tres modelos sobre el mismo escenario y devuelve las probabilidades de cada uno, de modo que pueden contrastarse en la interfaz. Los tres modelos se evalúan uno tras otro reutilizando el mismo vector de características, que solo se construye una vez.

  \item \textbf{Depuración de características} (\texttt{POST /api/lpmc/debug-features}). Devuelve el vector de características completo, con sus valores antes y después de aplicar el escalado estándar. Permite comprobar que las variables de ruta calculadas a partir de OSRM y OTP caen en el rango esperado por el dataset, y fue la herramienta principal para diagnosticar los problemas de unidades que se describen en la Sección~\ref{subsec:impl_lpmc_entrenamiento}.

  \item \textbf{Modelos disponibles} (\texttt{GET /api/lpmc/models}). Escanea el directorio \texttt{lpmc/models/} en busca de pares de ficheros \texttt{\{nombre\}\_lpmc.joblib} y \texttt{\{nombre\}\_lpmc\_scaler.joblib}, y devuelve la lista de modelos encontrados junto con el predeterminado (\texttt{LPMC\_MODEL\_VARIANT}). Gracias a este barrido, basta con dejar un modelo entrenado fuera del sistema en ese directorio para que la interfaz lo detecte en la siguiente petición, sin reiniciar ningún contenedor.
\end{itemize}

El nombre del modelo que llega desde el exterior se valida con la expresión regular \texttt{\textasciicircum[a-z0-9\_]\{1,32\}\$}, que solo admite letras minúsculas, dígitos y guiones bajos. La validación es necesaria por seguridad: ese nombre se concatena para formar la ruta del artefacto a cargar (\texttt{\{nombre\}\_lpmc.joblib}), y sin filtrarlo un valor manipulado como \texttt{../../} podría leer ficheros fuera del directorio de modelos. Por último, los modelos se cargan de forma diferida y quedan en una caché en memoria: la primera predicción con un modelo asume el coste de leerlo y deserializarlo desde disco, y las siguientes lo reutilizan sin volver a cargarlo.

\subsection{Resolución de rutas y portabilidad}
\label{subsec:impl_backend_docker}

El backend necesita localizar en disco los datos del feed GTFS y los artefactos de los modelos. En lugar de codificar rutas absolutas, que dependerían del equipo y del sistema operativo concretos, cada módulo deriva la ruta que necesita a partir de su propia ubicación con la biblioteca estándar \texttt{pathlib}. Por ejemplo, el módulo de inferencia reside en \texttt{movilidad-urbana-sim/backend/app/services/} y, para llegar a la raíz del repositorio, evalúa \texttt{Path(\_\_file\_\_).resolve().parents[4]}: cuatro niveles por encima de su propio fichero (\texttt{services} $\to$ \texttt{app} $\to$ \texttt{backend} $\to$ \texttt{movilidad-urbana-sim} $\to$ raíz). Desde ahí compone la ruta a \texttt{lpmc/models/}. El cálculo es correcto pero queda atado a esa profundidad concreta de carpetas; una alternativa más robusta sería ascender hasta encontrar un fichero marcador de la raíz del repositorio.

Esta forma de resolver rutas mantiene la compatibilidad entre el entorno de desarrollo (Windows) y el contenedor de ejecución (Linux), porque \texttt{pathlib} abstrae el separador de directorios de cada sistema y las rutas se calculan siempre relativas a la posición del código, no a un punto de montaje fijo.


% -------------------------------------------------------
\section{Implementación del frontend}
\label{sec:impl_frontend}
% -------------------------------------------------------

\subsection{Tecnologías y estructura}
\label{subsec:impl_frontend_tech}

El frontend utiliza las siguientes tecnologías principales:

\begin{itemize}
  \item \textbf{React~18} \cite{ReactDocs}: biblioteca de construcción de interfaces de usuario basada en componentes reutilizables y actualizaciones reactivas del árbol de elementos DOM.
  \item \textbf{Vite} \cite{ViteDocs}: herramienta de compilación que sirve los módulos de forma nativa durante el desarrollo, con actualización del módulo modificado sin recargar la página, y genera el paquete de producción optimizado.
  \item \textbf{TypeScript} \cite{TypeScriptDocs}: lenguaje que extiende JavaScript añadiendo tipos comprobados de forma estática. Resulta especialmente útil al modelar las respuestas del backend, como los itinerarios OTP con desglose de tramos, los feeds GTFS con múltiples entidades relacionadas o los vectores de probabilidades de los modelos.
  \item \textbf{Leaflet} \cite{LeafletDocs}, integrado mediante \textbf{React-Leaflet} \cite{ReactLeafletDocs}: biblioteca de mapas interactivos de código abierto. React-Leaflet expone cada primitiva cartográfica (marcadores, polilíneas, capas de teselas) como un componente React declarativo.
  \item \textbf{Lucide React} \cite{LucideReact}: colección de iconos SVG de código abierto, utilizada en los botones de la barra lateral y los controles del mapa. Al ser vectoriales, escalan sin pérdida de calidad.
  \item \textbf{TanStack Query} \cite{TanStackQuery}: librería de gestión del estado asíncrono para React que aporta caché de respuestas, seguimiento del ciclo de vida de cada petición (carga, error, éxito) y reintento automático.
\end{itemize}

La aplicación se estructura en dos componentes principales: \texttt{App} concentra el estado global y la lógica del cliente, y \texttt{MapView} recibe el estado necesario como \textit{props} (datos que el componente padre pasa al hijo de forma unidireccional) y es responsable exclusivamente del renderizado cartográfico. Esta separación se describe en la Sección~\ref{sec:arch_frontend}.

\subsection{Interfaz cartográfica y controles del mapa}
\label{subsec:impl_frontend_mapa}

La interfaz sigue el esquema de las aplicaciones cartográficas modernas: el mapa ocupa toda la ventana del navegador como fondo, mientras los controles y los paneles de navegación se superponen sobre él. Leaflet renderiza el mapa a partir de teselas de los proveedores seleccionados, sin descargar ningún mapa local; el extracto OSM que usa OSRM es solo para el cálculo de rutas en el servidor. Las rutas calculadas por OSRM cubren únicamente el territorio del extracto de Castilla-La Mancha, y las paradas de autobús disponibles pertenecen exclusivamente al feed GTFS de Toledo.

Los marcadores de origen y destino se implementan con iconos SVG propios embebidos en CSS, distintos de los marcadores predeterminados de Leaflet, para distinguir visualmente el punto de inicio del de llegada. Las paradas de autobús se representan como marcadores circulares (\texttt{CircleMarker}) en una capa de renderizado propia de Leaflet (\textit{pane} \texttt{stopsPane}, z-index~450 frente a los 400 de las polilíneas). Este orden de apilamiento asegura que las paradas sean siempre accesibles a la interacción del usuario aunque un trazado las atraviese.

Los controles de navegación del mapa se implementaron como componentes React, en lugar de los controles predeterminados de Leaflet, para seguir el mismo estilo visual que el resto de la interfaz. Se organizan en dos grupos fijos sobre el mapa:

\begin{itemize}
  \item \textbf{Esquina inferior derecha}: botones de acercar y alejar la vista, más un botón de recentrar que devuelve el mapa al centro de Toledo con el zoom inicial.
  \item \textbf{Esquina superior derecha}: tres botones de limpieza que eliminan, respectivamente, las rutas de navegación viaria (OSRM y OTP), la línea de bus activa en el panel Red GTFS, o ambas cosas junto con los puntos de origen y destino. Los botones se deshabilitan automáticamente cuando no hay datos que eliminar.
\end{itemize}

El zoom puede ajustarse en pasos de 0,25 en lugar de enteros, lo que aporta mayor precisión al encuadrar áreas urbanas concretas.

La barra lateral de navegación (\textit{rail}) de 64~px de ancho, siempre visible a la izquierda de la pantalla, da acceso a los seis paneles funcionales del simulador. Al pulsar un botón del \textit{rail}, el panel correspondiente se despliega a su derecha; pulsar de nuevo el mismo botón lo cierra. El panel activo al arrancar la aplicación es \textbf{Inicio}. Los demás paneles son \textbf{Rutas}, \textbf{Red GTFS}, \textbf{Predicción IA}, \textbf{Capas} y \textbf{Ajustes}, descritos en las secciones siguientes.

\begin{figure}[htbp]
  \centering
  \missingfigure{Captura general de la aplicación mostrando el mapa de Toledo con los cuatro modos de ruta dibujados simultáneamente, las paradas de bus visibles, el panel de Rutas abierto con el detalle de un itinerario OTP y el control de fecha/hora global en la esquina superior izquierda.}
  \caption{Vista general de la interfaz del simulador de movilidad urbana.}
  \label{fig:app_general}
\end{figure}

\subsection{Panel Inicio}
\label{subsec:impl_frontend_inicio}

El panel Inicio es la vista por defecto al cargar la aplicación. Presenta una descripción del proyecto y su contexto académico, una lista de las tecnologías utilizadas con sus versiones, y una tabla de créditos con autor, tutor, centro, dataset, feed GTFS y cartografía base. El objetivo es orientar al usuario que accede al simulador por primera vez sin que necesite consultar documentación externa. Cerrar el panel o cambiar a otro no afecta al estado del mapa ni a los cálculos en curso.

\subsection{Panel Rutas}
\label{subsec:impl_frontend_rutas}

El panel Rutas agrupa todas las funciones de navegación viaria y multimodal. Los puntos de origen y destino se establecen mediante el menú contextual (clic derecho sobre el mapa): la primera entrada muestra las coordenadas del punto pulsado con cinco decimales y las copia al portapapeles, con un cambio temporal del texto a <<¡Copiado!>> como confirmación; las dos entradas siguientes asignan el punto como origen o como destino. El mismo par de coordenadas puede editarse también directamente en los campos de texto del panel, pegando un valor «latitud, longitud» procedente de cualquier fuente; el campo valida el rango geográfico al confirmar con \textit{Intro} o al perder el foco, y restaura el valor previo si el formato no es válido.

Los botones de selección de modo de transporte no son mutuamente excluyentes: un clic normal selecciona en exclusiva el modo pulsado, mientras que un clic con \textit{Shift} añade o retira ese modo sin modificar el resto. Cuando hay varios modos activos, \texttt{MapView} renderiza una polilínea independiente por perfil: azul para conducción, verde para bicicleta, gris para desplazamiento a pie y naranja oscuro discontinuo para los tramos a pie del itinerario de transporte público. Las tres rutas OSRM se calculan siempre en paralelo, por lo que cambiar los modos activos no requiere ninguna nueva petición al servidor.

Al pulsar «Calcular rutas» se lanza la petición a \texttt{POST /api/osrm/routes} y, si el modo de transporte público está activo, también a \texttt{POST /api/otp/routes}. Una vez realizado ese primer cálculo, cualquier modificación posterior de los extremos O/D relanza ambas peticiones de forma automática. Este auto-recálculo se desactiva al limpiar las rutas y permanece inactivo hasta el siguiente cálculo manual, evitando peticiones antes de que el usuario haya definido el escenario.

Los resultados OSRM se presentan en una tabla con una fila por modo activo, mostrando distancia y duración estimada. El itinerario OTP se muestra en una tarjeta independiente con controles de paginación para navegar entre las alternativas devueltas. El encabezado de la tarjeta resume la hora de salida, la hora de llegada y la duración total; a continuación se despliega el detalle tramo a tramo. Los tramos a pie se resumen en una línea con distancia y duración. Cada tramo en autobús se encabeza con la etiqueta coloreada de la línea y las horas de salida y llegada, y bajo él se dibuja un diagrama vertical con todas las paradas: una barra del color de la línea, extremos como puntos rellenos, paradas intermedias como aros huecos y la hora de paso junto a cada parada. Al pulsar una parada del diagrama el mapa se desplaza hasta ella y se abre su ventana emergente. Cuando el itinerario incluye transbordos directos (cambio de línea en la misma parada sin tramo a pie entre ellos), el diagrama lo indica con una etiqueta de transbordo entre los dos bloques.

La etiqueta coloreada de cada línea en el panel debe coincidir con el color del catálogo GTFS. Sin embargo, OTP identifica las líneas con un identificador prefijado por el código del feed (por ejemplo, \texttt{1:50011}), que no coincide directamente con el \texttt{route\_id} del feed estático. La función auxiliar \texttt{resolveGtfsRoute()} resuelve la correspondencia en tres intentos sucesivos: por identificador exacto, por identificador sin el prefijo del feed y por nombre corto. Este proceso es necesario también para aplicar el contorno de contraste que se describe en la Sección~\ref{subsec:impl_frontend_red}.

La Figura~\ref{fig:otp_itinerario} muestra un itinerario multimodal con el desglose de tramos desplegado en la interfaz.

\begin{figure}[H]
  \centering
  \includegraphics[width=1.0\textwidth]{figs/ItinerarioMultimodalZoom.png}
  \caption{Itinerario multimodal OTP con desglose de tramos en la interfaz.}
  \label{fig:otp_itinerario}
\end{figure}

\begin{figure}[htbp]
  \centering
  \missingfigure{Captura del panel Rutas con la tabla de resultados OSRM para los tres modos activos y el itinerario OTP desplegado mostrando el diagrama de paradas de un tramo en autobús.}
  \caption{Panel Rutas con resultados OSRM e itinerario OTP.}
  \label{fig:panel_rutas}
\end{figure}

\subsection{Panel Red GTFS}
\label{subsec:impl_frontend_red}

El panel Red GTFS permite explorar la red de autobús urbano de Toledo sin necesidad de calcular un itinerario. Muestra las líneas disponibles, el trazado de la línea seleccionada sobre el mapa, la secuencia de paradas y los horarios del día activo.

El feed de Toledo contiene 56 sentidos (identificadores \texttt{route\_id} distintos) agrupados en 25 líneas por nombre corto (\texttt{route\_short\_name}). El panel organiza los contenidos en un acordeón: cada entrada del nivel superior representa una línea; las líneas con más de un sentido muestran un icono de despliegue y al expandirlas listan los sentidos disponibles con el sentido activo marcado por un borde lateral coloreado. El color de cada línea sigue una paleta fija de 26 colores definida en el cliente: los nombres cortos únicos se ordenan alfabéticamente y el i-ésimo recibe \texttt{LINE\_COLORS[i~mod~26]}, con lo que las 25 líneas del feed reciben colores únicos. Una vez determinado el color, la función \texttt{isLightColor()} calcula la luminancia perceptual con la fórmula $0{,}299R + 0{,}587G + 0{,}114B$; si el resultado supera 0,9 (líneas de color claro, como la L14 en blanco), se añade un contorno negro a las etiquetas coloreadas, a los marcadores de parada y al trazado del mapa para mantener la legibilidad sobre fondo blanco.

Al seleccionar una línea, el mapa dibuja su trazado como polilínea del color correspondiente y sus paradas como marcadores circulares rellenos con ese mismo color y un contorno de contraste. Los marcadores se renderizan en la capa \texttt{stopsPane} (z-index~450), por encima de las polilíneas (z-index~400), de modo que son siempre accesibles a la interacción aunque el trazado los cruce. Al pulsar una parada del listado lateral, la vista se desplaza hasta ella (\texttt{flyTo} a zoom~17) y, transcurridos 900~ms, se abre la ventana emergente de la parada. Las ventanas emergentes muestran nombre, código y las etiquetas coloreadas de las líneas que pasan por ella; al pulsar una etiqueta se selecciona ese sentido directamente en el panel.

Los horarios se presentan agrupados por hora: una columna con la hora y, a su derecha, los minutos de cada salida estilizados según su posición respecto a la hora activa del control global de fecha. Los minutos ya pasados aparecen en gris; los futuros, en azul; el próximo (menor tiempo restante), en negrita. Si no quedan salidas futuras ese día se muestra un aviso; si la línea no circula en la fecha seleccionada, el panel informa de ello sin tabla de horarios.

\begin{figure}[htbp]
  \centering
  \missingfigure{Captura del panel Red GTFS con el acordeón de líneas, una línea seleccionada con su diagrama de paradas y la tabla de horarios con la próxima salida destacada en negrita.}
  \caption{Panel Red GTFS con trazado, paradas y horarios de una línea.}
  \label{fig:panel_red}
\end{figure}

\subsection{Panel Predicción IA}
\label{subsec:impl_frontend_ia}

El panel Predicción IA, que muestra la Figura~\ref{fig:panel_ia}, permite obtener la probabilidad de cada modo de transporte para un viajero concreto en el escenario definido sobre el mapa. Su disposición sigue el flujo de los datos que recibe el modelo. En primer lugar, en modo de solo lectura, aparecen las coordenadas de origen y destino que se pasan a OSRM y OTP, de modo que el usuario verifica de un vistazo el escenario. Justo debajo, también en solo lectura, se muestran los tiempos y distancias que esos servicios calculan para cada modo de transporte (las variables de ruta que alimentan al modelo), de manera que el usuario ve exactamente con qué cifras se va a inferir. A continuación se indica la información del viaje (día de la semana y hora de salida), derivada del control global de fecha y hora; al no ser editable aquí, se evita cualquier inconsistencia entre el itinerario calculado y el perfil enviado al modelo.

\begin{figure}[htbp]
  \centering
  \missingfigure{Captura del panel Predicción IA mostrando, de arriba abajo: el par origen-destino en solo lectura, la tabla de tiempos y distancias por modo, el día y hora del viaje, los tres perfiles predefinidos, el formulario sociodemográfico y el resultado de la inferencia con la probabilidad de cada modo y el modo más probable destacado.}
  \caption{Panel Predicción IA con el formulario y el resultado de la inferencia.}
  \label{fig:panel_ia}
\end{figure}

Para facilitar la exploración del simulador, el panel ofrece tres perfiles sociodemográficos predefinidos que autocompletan todos los campos del formulario y ajustan además la fecha y la hora globales al escenario que representan:

\begin{itemize}
  \item \textbf{Commuter}: viaje al trabajo (HBW) en martes a las 8:15. Varón de 36~años con carnet de conducir y un vehículo en el hogar.
  \item \textbf{Estudiante}: viaje a la educación (HBE) en miércoles a las 7:45. Mujer de 21~años sin carnet de conducir ni vehículo en el hogar, con tarifa de transporte público reducida.
  \item \textbf{Familiar}: viaje de ocio (HBO) en sábado a las 11:30. Mujer de 44~años con carnet de conducir y dos vehículos en el hogar.
\end{itemize}

Al cargar la aplicación se aplica el perfil \textit{Commuter} por defecto, de modo que el panel arranca con un escenario completo y coherente. El indicador de perfil activo desaparece en cuanto el usuario modifica algún campo del formulario, eliminando la ambigüedad sobre qué valores se están usando. El resto de parámetros sociodemográficos (edad, género, carnet, vehículos en el hogar, motivo del viaje, tipo de combustible del vehículo, costes) se introducen en el formulario y son editables en cualquier momento. La inferencia se lanza con el botón «Inferir modo», que muestra junto a su rótulo el nombre del modelo activo. El resultado es un conjunto de probabilidades para los cuatro modos de transporte (a pie, bicicleta, transporte público y conducción), con el modo de mayor probabilidad destacado visualmente.

Tras la inferencia, bajo el botón aparecen las probabilidades de los cuatro modos en forma de barras y, solo entonces, se habilitan dos acciones secundarias para profundizar en el resultado. La primera, «Comparar modelos», llama a \texttt{POST /api/lpmc/compare} y presenta una tabla con las probabilidades de los tres modelos (XGBoost, Random Forest y red neuronal) para el mismo escenario, lo que permite ver hasta qué punto coinciden (Figura~\ref{fig:panel_comparacion}). La segunda, «Ver variables», llama a \texttt{POST /api/lpmc/debug-features} y abre una ventana modal con el vector de características completo, en sus valores originales y tras el escalado, útil para inspeccionar con qué datos exactos se ha inferido (Figura~\ref{fig:debug_features}).

\begin{figure}[htbp]
  \centering
  \missingfigure{Captura del panel de comparación de los tres modelos (XGBoost, Random Forest, DNN) mostrando las probabilidades de elección para los cuatro modos (a pie, bicicleta, transporte público y coche) ante el mismo par origen-destino y perfil de usuario. Mostrar un caso con clara concordancia entre modelos y el modo más probable destacado.}
  \caption{Panel de comparación de probabilidades de los tres modelos de elección modal.}
  \label{fig:panel_comparacion}
\end{figure}

\begin{figure}[htbp]
  \centering
  \missingfigure{Captura de la ventana modal «Ver variables» (debug-features) con la tabla del vector de características completo: cada variable con su valor original y su valor tras el escalado estándar, incluyendo las variables de ruta de OSRM/OTP y las sociodemográficas.}
  \caption{Ventana de depuración del vector de características que recibe el modelo.}
  \label{fig:debug_features}
\end{figure}

\subsection{Panel Capas}
\label{subsec:impl_frontend_capas}

El panel Capas permite seleccionar la capa de fondo del mapa entre seis opciones, cada una orientada a un tipo de análisis o preferencia visual:

\begin{itemize}
  \item \textbf{Color} (CartoDB Voyager): calles y nombres de lugar en tonos suaves. Seleccionada por defecto al iniciar la aplicación por ofrecer contexto urbano sin competir visualmente con las rutas superpuestas.
  \item \textbf{Blanco y negro} (CartoDB Positron): paleta completamente desaturada que minimiza el ruido visual y maximiza la legibilidad de los trazados.
  \item \textbf{OpenStreetMap}: mapa estándar de la comunidad OSM, con mayor detalle de nombres e infraestructuras.
  \item \textbf{Topográfico} (OpenTopoMap): curvas de nivel y sombreado de relieve, especialmente relevante para trayectos en bicicleta o a pie por zonas con desnivel.
  \item \textbf{Satélite} (Esri World Imagery): ortofoto aérea global con resolución variable según la región.
  \item \textbf{PNOA} (Plan Nacional de Ortofotografía Aérea, IGN España): ortofoto servida mediante el protocolo WMTS por el Instituto Geográfico Nacional; cubre el territorio español con resolución de 25~cm en zonas urbanas y no requiere token de autenticación.
\end{itemize}

La selección de capa persiste mientras el usuario navega entre los demás paneles. Un componente auxiliar (\texttt{BasemapZoomSnapper}) detecta el cambio de proveedor y redondea el nivel de zoom al entero más próximo si es fraccionario, evitando errores de carga con los servidores de teselas que no sirven niveles no enteros.

\subsection{Panel Ajustes}
\label{subsec:impl_frontend_ajustes}

El panel Ajustes agrupa tres secciones de configuración global. La primera sincroniza la visibilidad de las paradas de autobús en el mapa con el interruptor del panel Red GTFS, permitiendo activarlas o desactivarlas desde cualquiera de los dos puntos de la interfaz sin inconsistencias. La segunda aloja el selector de modelo activo de inferencia: al abrirse el panel por primera vez, la interfaz consulta \texttt{GET /api/lpmc/models} para obtener la lista de modelos disponibles e inicializa el selector con el valor de \texttt{default\_variant} que devuelve la API, que a su vez refleja el contenido de la variable de entorno \texttt{LPMC\_MODEL\_VARIANT}. Cambiar la selección no dispara ninguna petición inmediata; el identificador del modelo elegido se adjunta como campo del cuerpo de la siguiente petición a \texttt{POST /api/lpmc/predict}. La tercera sección contiene instrucciones colapsables para añadir modelos propios: convenio de nombres (\texttt{\{nombre\}\_lpmc.joblib} y \texttt{\{nombre\}\_lpmc\_scaler.joblib} en \texttt{lpmc/models/}), formatos admitidos (cualquier objeto sklearn con interfaz \texttt{predict\_proba}, o una red PyTorch cargable mediante \texttt{TorchModalWrapper}) y cómo fijar el modelo predeterminado. Cualquier par de ficheros que siga el convenio es detectado automáticamente por \texttt{GET /api/lpmc/models} sin modificar el código ni reiniciar los contenedores.

La Figura~\ref{fig:panel_ajustes} muestra el panel con los tres modelos incluidos en el repositorio disponibles en el desplegable.

\begin{figure}[htbp]
  \centering
  \missingfigure{Captura del panel Ajustes mostrando: (1) checkbox de paradas de bus, (2) selector de modelo activo con los tres modelos disponibles (XGBoost, Random Forest, DNN) y la etiqueta ``Activo'' junto al seleccionado, (3) bloque de instrucciones colapsable para añadir modelos personalizados.}
  \caption{Panel de ajustes con selector de modelo activo y documentación de extensión.}
  \label{fig:panel_ajustes}
\end{figure}

\subsection{Gestión del estado con TanStack Query}
\label{subsec:impl_frontend_query}

Las peticiones al backend se gestionan con TanStack Query (\texttt{@tanstack/react-query}). Los datos GTFS estáticos (paradas, lista de rutas) se cargan con \texttt{useQuery} al arrancar la aplicación y se mantienen en caché durante la sesión. Las peticiones de rutas OSRM, itinerarios OTP e inferencia LPMC se declaran como mutaciones (\texttt{useMutation}), que se disparan manualmente al pulsar los controles de cálculo en la interfaz.

TanStack Query gestiona de forma automática los estados de carga y error, la caché y el reintento con retroceso exponencial (cada reintento espera más que el anterior). Las tres mutaciones de inferencia LPMC (\texttt{/predict}, \texttt{/compare}, \texttt{/debug-features}) se declaran por separado, de modo que la predicción con el modelo activo y la comparación de los tres modelos se lanzan de forma independiente, sin afectar al estado del resto de la interfaz.

% -------------------------------------------------------
\section{Modelos de elección modal}
\label{sec:impl_lpmc}
% -------------------------------------------------------

Esta sección describe el proceso de entrenamiento, validación y despliegue en producción de los tres modelos de elección modal que integra el simulador. La arquitectura del pipeline de inferencia se describe en la sección~\ref{sec:arch_pipeline_inferencia}; aquí se detalla la implementación concreta: el dataset utilizado, las decisiones de preprocesado, la configuración de los modelos y los resultados obtenidos.

\subsection{Dataset LPMC y variables de entrada}
\label{subsec:impl_lpmc_dataset}

El dataset London Passenger Mode Choice (LPMC) recoge registros de viajes reales observados en el área metropolitana de Londres entre 2012 y 2015. Procede de la \textit{London Travel Demand Survey} (LTDS), la encuesta de movilidad de Transport for London, sobre la que Hillel et al.\ construyeron el dataset que aquí se utiliza: a cada viaje observado le añadieron las variables de nivel de servicio, esto es, los tiempos y costes de cada modo de transporte alternativo (a pie, bicicleta, transporte público y coche) para ese mismo par origen-destino, calculados aunque el viajero no eligiera ese modo. Se citan dos publicaciones porque corresponden a aportaciones distintas: el artículo original que presenta el dataset \cite{Hillel2018LPMC} y su descripción técnica detallada \cite{CSLPMC2019}. El dataset es, además, el utilizado en los experimentos de referencia sobre los que se apoya este trabajo \cite{MartinBaos2023Thesis,MartinBaos2023TRC}. La Tabla~\ref{tab:lpmc_dataset_stats} resume sus características principales.

\begin{table}[htbp]
  \caption{Características generales del dataset LPMC.}
  \label{tab:lpmc_dataset_stats}
  \centering
  \begin{tabular}{lc}
    \toprule
    \textbf{Característica} & \textbf{Valor} \\
    \midrule
    Registros totales        & 81.086 \\
    Hogares únicos           & 17.616 \\
    Modos de viaje           & 4 (\textit{walk}, \textit{cycle}, \textit{pt}, \textit{drive}) \\
    Oleadas de encuesta      & 3 (\texttt{survey\_year} 1, 2, 3) \\
    Años de viaje cubiertos  & 2012--2015 \\
    Variables originales     & 34 \\
    \bottomrule
  \end{tabular}
\end{table}

El interés metodológico del dataset para este TFM no reside en el caso geográfico concreto (Londres), sino en que sus variables de nivel de servicio son justamente las que OSRM y OTP calculan para cualquier par origen-destino (tiempos por modo, transbordos, costes). Contiene también las variables sociodemográficas que el simulador pide al usuario, tiene un tamaño suficiente para entrenar modelos de aprendizaje automático con validación cruzada robusta, y ya ha sido validado para esta misma tarea en investigaciones previas \cite{MartinBaos2023TRC}.

Las variables de entrada al modelo se dividen en dos grupos. El primero comprende las variables de ruta, derivadas automáticamente de OSRM y OTP para cada par origen-destino consultado en el simulador. El segundo comprende las variables sociodemográficas y contextuales, introducidas por el usuario a través del panel lateral de la interfaz o autocompletadas a partir de uno de los perfiles predefinidos. El listado completo de variables, con su tipo de dato y su origen, se recoge en la Tabla~\ref{tab:lpmc_features} del Anexo~\ref{anx:experimentos}.

\subsection{Preprocesado de datos}
\label{subsec:impl_lpmc_preprocesado}

El preprocesado del dataset transforma los datos en bruto en el conjunto de variables que reciben los modelos. Se aplica en cuatro pasos.

\begin{enumerate}
  \item \textbf{Eliminación de columnas sin valor predictivo.} Se descartan doce columnas. Tres son identificadores (\texttt{trip\_id}, \texttt{person\_n}, \texttt{trip\_n}), que no aportan señal. Tres son variables de fecha de grano fino (\texttt{travel\_year}, \texttt{travel\_month}, \texttt{travel\_date}), cuya información ya recogen \texttt{day\_of\_week} y \texttt{start\_time\_linear}. Una es un factor de política exógeno que no está disponible en inferencia (\texttt{bus\_scale}). Las cinco restantes son variables derivadas redundantes, que se obtienen sumando otras ya presentes (\texttt{dur\_pt\_total}, \texttt{dur\_pt\_int\_total}, \texttt{cost\_driving\_fuel}, \texttt{cost\_driving\_con\_charge}, \texttt{driving\_traffic\_percent}).

  \item \textbf{Codificación de variables categóricas.} Las dos variables categóricas se transforman a representación \textit{one-hot}, esto es, una columna binaria por categoría. El motivo del viaje (\texttt{purpose}) genera cinco columnas (\textit{B, HBE, HBO, HBW, NHBO}). El tipo de combustible (\texttt{fueltype}) se consolida primero de seis categorías a cuatro (\textit{Petrol, Diesel, Hybrid, Average}), unificando las variantes de furgoneta ligera (\texttt{Petrol\_LGV}, \texttt{Diesel\_LGV}) con sus equivalentes de turismo, y después se codifica con \textit{one-hot}. La variable objetivo \texttt{travel\_mode} se convierte a entero (\textit{walk}~=~0, \textit{cycle}~=~1, \textit{pt}~=~2, \textit{drive}~=~3).

  \item \textbf{División temporal de los datos.} El conjunto se separa por la oleada de encuesta (\texttt{survey\_year}), no de forma aleatoria. Las oleadas 1 y 2 forman el conjunto de entrenamiento y la oleada 3, la más reciente, el conjunto de prueba. Con una división aleatoria, viajes de un mismo periodo caerían a la vez en entrenamiento y prueba, y el modelo podría aprender tendencias del periodo de prueba en lugar de generalizar hacia el futuro. La separación temporal reproduce la condición real de uso del simulador, que predice sobre escenarios posteriores a los datos con los que se entrenó.

  \item \textbf{Normalización de variables continuas.} Las dieciséis variables continuas (distancias, duraciones por modo, transbordos, edad, coches del hogar, hora, día y costes) se normalizan con \texttt{StandardScaler} de scikit-learn, que resta la media y divide por la desviación típica de cada variable para dejarla con media~0 y desviación típica~1. Las variables binarias y las columnas \textit{one-hot} no se escalan. Los modelos de árbol (XGBoost y Random Forest) son insensibles al escalado, pero se aplica el mismo preprocesado a los tres modelos para que el pipeline de inferencia del backend sea idéntico con cualquiera de ellos.
\end{enumerate}

El escalador se ajusta siempre solo con datos de entrenamiento y nunca con los de evaluación. En la validación cruzada se reajusta sobre el subconjunto de entrenamiento de cada \textit{fold}; para el modelo final, sobre todo el conjunto de entrenamiento. De este modo, la media y la desviación típica que usa el escalador no incorporan información del conjunto de validación, lo que evitaría una fuga de datos que daría una estimación de error demasiado optimista.

\subsection{Ajuste de hiperparámetros}
\label{subsec:impl_lpmc_hiperparametros}

Los hiperparámetros son los ajustes que gobiernan cómo aprende cada modelo (por ejemplo, el número de árboles o la tasa de aprendizaje) y que no se estiman durante el entrenamiento, sino que hay que fijar de antemano. Para los tres modelos se realizó una búsqueda propia de hiperparámetros con una metodología común: se trabajó sobre una muestra del conjunto de entrenamiento (para acotar el tiempo de cómputo), se evaluó cada configuración con validación cruzada agrupada por hogar y se tomó como criterio la entropía cruzada media en validación, equivalente a maximizar el GMPCA descrito en la Sección~\ref{subsec:impl_lpmc_entrenamiento}. Esta búsqueda nos da una configuración de partida razonable para cada modelo. En el caso particular de XGBoost, además, se dispone de la búsqueda exhaustiva publicada en la investigación de referencia \cite{MartinBaos2023Thesis}, mucho más costosa y ya validada, que se adopta para el modelo final por encima de la nuestra. Los apartados siguientes detallan cada modelo.

\subsubsection*{Random Forest}

La búsqueda para el Random Forest se hizo con optimización Bayesiana (HyperOpt con el estimador de Parzen estructurado en árbol, TPE), 40 evaluaciones sobre el 25\,\% del conjunto de entrenamiento. El TPE construye un modelo probabilístico de la relación entre los hiperparámetros y el error, y dirige cada nueva evaluación hacia las regiones del espacio que mejores resultados han dado, de modo que converge antes que una búsqueda en rejilla (que probaría todas las combinaciones) o que una búsqueda aleatoria (que muestrea sin aprender de las pruebas anteriores). El espacio explorado abarcó el número de árboles, la profundidad máxima, los tamaños mínimos de nodo y hoja, y el número de variables candidatas por división. La mejor configuración encontrada (Tabla~\ref{tab:rf_params}) se adoptó como modelo final.

\begin{table}[htbp]
  \caption{Hiperparámetros del Random Forest hallados por la búsqueda propia (HyperOpt/TPE, 40 evaluaciones).}
  \label{tab:rf_params}
  \centering
  \begin{tabular}{ll}
    \toprule
    \textbf{Parámetro} & \textbf{Valor} \\
    \midrule
    \texttt{n\_estimators}      & 600 \\
    \texttt{max\_depth}         & 20 \\
    \texttt{min\_samples\_split} & 9 \\
    \texttt{min\_samples\_leaf} & 5 \\
    \texttt{max\_features}      & \texttt{sqrt} \\
    \bottomrule
  \end{tabular}
\end{table}

Los 600 árboles del bosque se entrenan cada uno sobre una muestra distinta de los datos y de las variables (\textit{bagging}); promediar sus predicciones reduce la varianza del conjunto. La profundidad máxima de 20 niveles, junto con los mínimos de 9 muestras para dividir un nodo y 5 para formar una hoja, impide que los árboles crezcan hasta memorizar casos concretos, lo que limita el sobreajuste. El parámetro \texttt{max\_features=sqrt} indica que en cada división solo se considera la raíz cuadrada del número total de variables; es el valor que Breiman recomienda para clasificación \cite{Breiman2001}, ya que decorrelaciona los árboles y mejora la capacidad de generalización del conjunto.

\subsubsection*{XGBoost}

Para XGBoost se ejecutó primero una búsqueda propia equivalente a la del Random Forest (HyperOpt/TPE, 80 evaluaciones sobre el 25\,\% del conjunto de entrenamiento), que sirvió como base de comparación. No obstante, la investigación de referencia \cite{MartinBaos2023Thesis} publica una búsqueda mucho más exhaustiva sobre este mismo dataset, con 1.000 evaluaciones y mayor presupuesto de cómputo, cuyos hiperparámetros (Tabla~\ref{tab:xgb_params}) producían mejores resultados que los de nuestra búsqueda reducida. Por ese motivo se adoptaron directamente esos valores para el modelo final, en lugar de los de la búsqueda propia.

\begin{table}[htbp]
  \caption{Hiperparámetros de XGBoost adoptados de la búsqueda de referencia (HyperOpt/TPE, 1.000 evaluaciones).}
  \label{tab:xgb_params}
  \centering
  \begin{tabular}{ll}
    \toprule
    \textbf{Parámetro} & \textbf{Valor} \\
    \midrule
    \texttt{n\_estimators}      & 4.809 \\
    \texttt{learning\_rate}     & 0,01 \\
    \texttt{max\_depth}         & 4 \\
    \texttt{min\_child\_weight} & 35 \\
    \texttt{gamma}              & 3,93 \\
    \texttt{max\_delta\_step}   & 9 \\
    \texttt{subsample}          & 0,767 \\
    \texttt{colsample\_bytree}  & 0,934 \\
    \texttt{colsample\_bylevel} & 0,651 \\
    \texttt{reg\_alpha} (L1)    & 0,048 \\
    \texttt{reg\_lambda} (L2)   & 0,037 \\
    \bottomrule
  \end{tabular}
\end{table}

La combinación de una tasa de aprendizaje baja (\texttt{learning\_rate=0,01}) con un número elevado de árboles (4.809) es característica de un modelo robusto: con pasos pequeños el modelo necesita más iteraciones para converger, pero el resultado tiene menor varianza. La penalización \texttt{gamma=3,93} y el peso mínimo de hoja \texttt{min\_child\_weight=35} actúan como regularización estructural, ya que evitan divisiones que no aporten una ganancia suficiente y reducen el sobreajuste en las clases minoritarias, como la bicicleta.

\subsubsection*{Red neuronal profunda}

Como cada evaluación de la red exige un ciclo completo de entrenamiento, en lugar de la búsqueda Bayesiana se realizó una búsqueda aleatoria acotada: se muestrearon 12 configuraciones del espacio de hiperparámetros (tamaños de las capas, tasa de aprendizaje, \textit{dropout}, regularización y tamaño de lote), cada una evaluada con validación cruzada agrupada por hogar sobre la mitad del conjunto de entrenamiento. La mejor configuración (Tabla~\ref{tab:dnn_params}) se adoptó para el modelo final.

\begin{table}[htbp]
  \caption{Configuración de la red neuronal profunda hallada por la búsqueda propia.}
  \label{tab:dnn_params}
  \centering
  \begin{tabular}{ll}
    \toprule
    \textbf{Hiperparámetro} & \textbf{Valor} \\
    \midrule
    Neuronas por capa oculta & 256, 128, 64 \\
    \textit{Dropout} (capas 1 y 2) & 0,4 y 0,3 \\
    Optimizador              & Adam \\
    Tasa de aprendizaje      & 0,003 \\
    Regularización L2 (\texttt{weight\_decay}) & 0,0001 \\
    Tamaño de lote           & 512 \\
    Función de pérdida       & Entropía cruzada (\textit{label smoothing} 0,1) \\
    \bottomrule
  \end{tabular}
\end{table}

La red, representada en la Figura~\ref{fig:dnn_arch}, transforma el vector de entrada a través de tres capas ocultas de 256, 128 y 64 neuronas y una capa final de 4 salidas, una por modo de transporte. Cada capa oculta encadena cuatro operaciones. Primero, una capa lineal (\textit{Linear}) combina sus entradas mediante una suma ponderada. Después, la normalización por lotes (\textit{BatchNorm}) reescala esas sumas para que tengan media y escala estables dentro de cada lote de entrenamiento, lo que evita que los valores crezcan o se desvanezcan al atravesar las capas y hace que el entrenamiento sea más estable. A continuación, la función de activación ReLU introduce la no linealidad que permite a la red aprender relaciones complejas, dejando pasar los valores positivos y anulando los negativos. Por último, en las dos primeras capas, el \textit{dropout} desactiva al azar una fracción de las neuronas (40\,\% y 30\,\%, respectivamente) en cada paso de entrenamiento, lo que obliga a la red a no depender de neuronas concretas y reduce el sobreajuste.

El orden de estas operaciones importa: la normalización se aplica después de la capa lineal y antes de la activación, sobre las sumas ponderadas y no sobre la salida de la capa anterior, lo que mantiene la entrada a la ReLU en un rango estable a lo largo del entrenamiento. La red se entrena con el optimizador Adam, que adapta el tamaño del paso de actualización para cada parámetro, y la función de pérdida es la entropía cruzada con suavizado de etiquetas (\textit{label smoothing} de 0,1): en lugar de exigir al modelo una probabilidad de 1 para la clase correcta y 0 para el resto, se le pide repartir una pequeña parte de la probabilidad entre todas las clases, lo que evita que se vuelva excesivamente confiado y mejora el ajuste de sus probabilidades. El entrenamiento incorpora además tres mecanismos de control: reducción de la tasa de aprendizaje a la mitad cuando la pérdida de validación deja de mejorar (\texttt{ReduceLROnPlateau}), recorte de la magnitud de los gradientes para evitar actualizaciones bruscas (\texttt{clip\_grad\_norm=1,0}) y parada anticipada, que detiene el entrenamiento y recupera el mejor estado si la pérdida de validación no mejora en 10 épocas seguidas. La semilla aleatoria se fijó a 481516 en PyTorch y NumPy para que el entrenamiento sea reproducible.

\begin{figure}[htbp]
  \centering
  \missingfigure{Diagrama de la arquitectura de la red neuronal: vector de entrada de n características → bloque 1 [Linear(256) → BatchNorm → ReLU → Dropout 0,4] → bloque 2 [Linear(128) → BatchNorm → ReLU → Dropout 0,3] → bloque 3 [Linear(64) → BatchNorm → ReLU] → capa de salida Linear(4) → softmax → probabilidades de walk, cycle, pt, drive.}
  \caption{Arquitectura de la red neuronal profunda de elección modal.}
  \label{fig:dnn_arch}
\end{figure}

\subsection{Entrenamiento y evaluación}
\label{subsec:impl_lpmc_entrenamiento}

Los tres modelos se entrenaron y validaron con validación cruzada de 5~\textit{folds} agrupada por hogar (\texttt{GroupKFold} de scikit-learn, con \texttt{household\_id} como variable de agrupación). Esta estrategia obliga a que todos los viajes de un mismo hogar queden en el mismo \textit{fold}, nunca repartidos entre entrenamiento y validación. De no ser así se produciría una fuga de datos (\textit{data leakage}), porque los viajes de un mismo hogar comparten las variables del hogar (como el número de coches o la posesión de carnet) y patrones de movilidad muy parecidos: el modelo podría reconocer al hogar y memorizar su comportamiento en lugar de aprender a generalizar, lo que daría una estimación de error demasiado optimista. Al agrupar por hogar, la validación se aproxima a un escenario \textit{leave-one-household-out}, en el que cada hogar se evalúa con un modelo que no ha visto ninguno de sus viajes, que es justamente la situación real cuando el simulador predice para un viajero nuevo. Las métricas de la Tabla~\ref{tab:lpmc_train_results} son la media de los cinco \textit{folds}.

La evaluación combina dos métricas. La primera es la exactitud (\textit{accuracy}), la proporción de viajes cuyo modo de transporte predicho (el de mayor probabilidad) coincide con el modo real:
\begin{equation}
  \label{eq:accuracy}
  \text{Accuracy} = \frac{1}{N}\sum_{i=1}^{N} \mathbb{1}\!\left[\hat{y}_i = y_i\right]
\end{equation}
donde $N$ es el número de viajes evaluados, $y_i$ el modo real del viaje $i$, $\hat{y}_i$ el modo predicho y $\mathbb{1}[\cdot]$ una función indicadora que vale 1 cuando ambos coinciden y 0 en caso contrario. La exactitud solo atiende al modo más probable e ignora con cuánta seguridad lo predice el modelo. Para medir también la calidad de las probabilidades se parte de la entropía cruzada media $H$:
\begin{equation}
  \label{eq:crossentropy}
  H = -\frac{1}{N}\sum_{i=1}^{N} \ln \hat{p}(y_i)
\end{equation}
donde $\hat{p}(y_i)$ es la probabilidad que el modelo asigna al modo realmente elegido en el viaje $i$. La entropía cruzada penaliza asignar probabilidades bajas a la clase correcta, tanto más cuanto más confiada y equivocada sea la predicción. A partir de ella se define el GMPCA (\textit{Geometric Mean Probability of Correct Alternative}):
\begin{equation}
  \label{eq:gmpca}
  \text{GMPCA} = \exp(-H) = \exp\!\left(\frac{1}{N}\sum_{i=1}^{N} \ln \hat{p}(y_i)\right)
\end{equation}
esto es, la media geométrica de las probabilidades que el modelo asigna al modo correcto. Un clasificador que repartiera probabilidad uniforme entre los cuatro modos obtendría GMPCA~=~0,25, y un clasificador perfecto, GMPCA~=~1,0. Frente a la exactitud, el GMPCA recompensa que el modelo asigne probabilidades altas al modo correcto y es la métrica más recomendada para la comparación en la literatura de elección modal con aprendizaje automático \cite{MartinBaos2023TRC}.

\begin{table}[htbp]
  \caption{Resultados en entrenamiento con validación cruzada de 5-folds agrupada por hogar (dataset LPMC).}
  \label{tab:lpmc_train_results}
  \centering
  \begin{tabular}{lcc}
    \toprule
    \textbf{Modelo} & \textbf{Accuracy (\%)} & \textbf{GMPCA (\%)} \\
    \midrule
    Random Forest & 74,95 & 51,51 \\
    XGBoost       & 75,48 & 52,54 \\
    DNN           & 75,17 & 51,31 \\
    \bottomrule
  \end{tabular}
\end{table}

\begin{table}[htbp]
  \caption{Resultados en el conjunto de prueba independiente (dataset LPMC).}
  \label{tab:lpmc_test_results}
  \centering
  \begin{tabular}{lcc}
    \toprule
    \textbf{Modelo} & \textbf{Accuracy (\%)} & \textbf{GMPCA (\%)} \\
    \midrule
    Random Forest & 74,03 & 50,60 \\
    XGBoost       & 74,41 & 51,63 \\
    DNN           & 74,35 & 50,36 \\
    \bottomrule
  \end{tabular}
\end{table}

Las Tablas~\ref{tab:lpmc_train_results} y~\ref{tab:lpmc_test_results} muestran que los tres modelos alcanzan un rendimiento muy similar, con XGBoost ligeramente por delante en ambas métricas. La diferencia entre la validación cruzada y la prueba se mantiene por debajo de un punto porcentual en todos los modelos, señal de que la división temporal no introduce diferencias apreciables entre las oleadas del dataset. El GMPCA, en torno al 50--52\,\%, duplica aproximadamente al de un clasificador aleatorio sobre cuatro clases (25\,\%), en línea con los resultados de investigaciones previas sobre el mismo dataset \cite{MartinBaos2023TRC}. XGBoost obtiene el mayor valor de GMPCA en prueba (51,63\,\%), una ligera ventaja tanto en exactitud como en la calidad de las probabilidades que asigna.

La elección de XGBoost como modelo principal del simulador no responde solo a esta ventaja numérica, que es estrecha, sino sobre todo a su comportamiento en la práctica. Como se explica a continuación, es el modelo que responde de forma más fiable y representativa al conocimiento experto que el backend inyecta ante escenarios extremos, lo que en un simulador interactivo importa más que unas décimas de exactitud.

Estos resultados deben enmarcarse en una limitación de fondo de los modelos de elección modal, tanto los clásicos de elección discreta como los de aprendizaje automático. Todos maximizan una utilidad y asumen un viajero racional, pero las decisiones reales incorporan un componente de ruido irreducible: no es realista esperar una probabilidad exactamente nula para un modo, de modo que una probabilidad residual del orden del 1\,\% no debe interpretarse como un error del modelo. Por otra parte, todo modelo de aprendizaje automático pierde fiabilidad al extrapolar fuera del rango de datos con el que se entrenó. El dataset LPMC, por ejemplo, contiene pocos trayectos muy cortos, y en ellos el modelo tiende a preferir el coche aunque a pie se tarde menos. En esas regiones poco cubiertas por los datos resulta útil inyectar conocimiento experto que corrija el comportamiento sin tocar el modelo: la penalización del transporte público cuando no hay servicio real y el recargo por aparcamiento que el backend suma al coche en trayectos muy cortos (Sección~\ref{subsec:impl_backend_penalizacion}) son ejemplos de este enfoque, que ajusta las entradas allí donde los datos no bastan.

El comportamiento de los tres modelos ante la penalización de itinerarios sin transporte público (Sección~\ref{subsec:impl_backend_penalizacion}) resultó revelador. Esa penalización lleva las características de transporte público a valores extremos, muy alejados de los vistos en entrenamiento. XGBoost responde como se espera y reduce la probabilidad del transporte público a un valor casi nulo. El Random Forest, en cambio, está acotado por construcción: su predicción es el promedio de las distribuciones de clase almacenadas en las hojas, por lo que la probabilidad del transporte público desciende pero no llega a anularse, quedándose en un valor residual que el usuario percibe como una opción todavía algo probable. La red neuronal no tiene esa cota: ante entradas tan alejadas de la distribución de entrenamiento puede extrapolar de forma descontrolada y concentrar casi toda la probabilidad en un único modo de transporte. La literatura atribuye a las redes neuronales una capacidad de extrapolación más flexible que la de los modelos de árbol \cite{Goodfellow2016}, pero esa flexibilidad no asegura que la extrapolación sea acertada en un caso concreto. Por todo ello, tanto por su rendimiento como por su respuesta más predecible ante los escenarios límite del simulador, se escogió XGBoost como modelo principal, sin necesidad de corregir su salida después de la inferencia.
